\documentclass[10pt]{article}

\usepackage[utf8]{inputenc}

\usepackage[T1]{fontenc}

\usepackage{amsmath}

\usepackage{amsfonts}

\usepackage{amssymb}

\usepackage[version=4]{mhchem}

\usepackage{stmaryrd}



\begin{document}

Больцмана, $T$ - абсолютная температура, $\mu$ - химич. потенциал, определяемый из условия нормировки: сумма всех $\bar{n}_{i}$ равна полному числу частиц в системе $\sum_{i} \bar{n}_{i}=N$.



Б. -Э. р.- функция распределения идеального газа из частиц, подчиняющихся квантовой Бозе - Эйнштейна статистике. В отсутствие внешнего поля Б.-Э. р. частиц идеального газа по абсолютной величине импульса $\boldsymbol{p}$ в узком слое $d p$ в объеме $V$ имеет вид



\[

d N_{p}=\left[\exp \left(\frac{p^{2}}{2 m}-\mu\right) / k T-1\right]^{-1} \frac{g V p^{2}}{2 \pi^{2} \hbar^{3}},

\]



где $m$ - масса частицы, $s-$ ее спин, $g=2 s+1-$ фактор вырождения квантовых состояний.



Из условия нормировки следует, что $\mu$ при $N \rightarrow \infty$ конечен лишь при температуре



\[

T>T_{0}=3,31 \frac{\hbar^{2}}{m k}\left(\frac{N}{V}\right)^{2 / 3},

\]



$T_{0}$ - температура вырождения.



При $T<T_{0}$ химич. потенциал $\mu$ обращается в нуль в термодинамич. пределе (когда $V \rightarrow \infty, N \rightarrow \infty, N / V-$ конечно) и в состоянии $\boldsymbol{p}=0$ (для покояшегося газа) оказывается макроскопически большое число частиц $N_{0}$ :



\[

N_{0}=N\left(1-\left(T / T_{0}\right)^{3 / 2}\right)

\]



При температуре $T_{0}$ длина волны де Бройля становится порядка среднего расстояния между частицами. При $T<T_{0}$ Б.-Э. р. имеет вид



\[

d N_{p}=\left[\exp \left(p^{2} / 2 m k T\right)-1\right]^{-1} \frac{g V p^{2} d p}{2 \pi^{2} \hbar^{3}} .

\]



Б. -Э. р. применимо к фотонам и квазичастицам, напр. фононам в кристаллич. решетке и в Не II. Для газов малой плотности, когда можно пренебречь квантовым эффектом, Б.-Э. р. переходит в Больцмана распределение.



Jum.: [1] B ose S. N., "Z. Phys.», 1924, Bd 27, S. 384-93; [2] Эй нштейн А., Собрание научных трудов, т. 3, М., 1966, с. 481-502. Д. Н. Зубарев.



БОЗЕ-ЭЙНШТЕЙНА СТАТИСТИКА, бозе-статис т и Ка,- квантовая статистика, применяемая к системам тождественных частиц с нулевым или целым спином (в единицах ћ). Предложена Ш. Бозе (Sh. Bose, 1924) для фотонов и обобщена А. Эйнштейном (A. Einstein, 1924) для одноатомного идеального газа. Характерная особенность Б. -Э. с. заключается в том, что в одном и том же квантовом состоянии может находиться любое число частиц.



B. Паули (W. Pauli, 1940) доказал, что тип квантовой статистики однозначно связан со значением спина частиц, так что совокупности частиц с нулевым или целым спином (ядра с четным числом нуклонов, фотоны, $\pi$-мезоны и другие так наз. бозоны) подчиняются Б.-Э. с., а системы частиц с полуцелым спином (электроны, нуклоны, ядра с нечетным числом нуклонов и другие так наз. фермионы) подчиняются Ферми - Дирака статистике.



При квантовомеханич. описании состояние системы описывается волновой функцией, к-рая в случае тождественных частиц либо симметрична по отношению к перестановкам любой пары частиц (для частиц с целым или нулевым спином), либо антисимметрична (для частиц с полуцелым спином). Для системы частиц, подчиняющихся Б.-Э.с., состояния описываются симметричными функциями, что является другой, эквивалентной формулировкой Б.-Э.с. Подобные системы называются бозе-системам и, напр. бозе - газ - газ из частиц, подчиняющихся Б.-Э. с.



Для идеального бозе-газа среднее число частиц $\bar{n}_{i}$ в состоянии $i$ определяется Бозе - Эйнштейна распределением. Формула для $\bar{n}_{i}$ следует из Гиббса распределения для идеального квантового газа с уровнями энергии



\[

E_{n}=\sum_{i} E_{i} n_{i}

\]



где $n_{i}$ согласно Б.-Э. с. могут принимать лишь значения $0,1,2, \ldots$ Распределение Бозе-Эйнштейна можно получить другим методом, если рассматривать статистически равновесное состояние квантового идеального газа как наиболее вероятное состояние и с помощью комбинаторики, учитывая неразличимость частиц, найти термодинамич. вероятность (с т а т и с т и ч е с к и й в е с) такого состояния, то есть число способов реализации данного квантового состояния газа с заданной энергией $E$ и числом частиц $N$. Для больших систем, когда $N$ велико, уровни энергии расположены очень плотно и стремятся к непрерывному распределению при стремлении числа частиц и объема к бесконечности. Пусть уровни сгруппированы по малым ячейкам, содержащим $G_{i}$ уровней в ячейке, число $G_{i}$ предполагается очень большим. Каждой $i$-й ячейке соответствует средняя энергия $E_{i}$ и число частиц $N_{i}$. Состояние системы определяется набором чисел $N_{i}$, где $N_{i}$ - сумма $n_{i}$ по уровням ячейки. Для Б.-Э. с. атомы предполагаются неразличимыми и в каждой ячейке может находиться произвольное число частиц. Поэтому статистич. вес $W\left(N_{i}\right)$ равен числу различных распределений частиц по ячейкам:



\[

W\left(N_{i}\right)=\prod_{i} \frac{\left(G_{i}+N_{i}-1\right) !}{N_{i} !\left(G_{i}-1\right) !},

\]



он определяет вероятность распределения частиц по ячейкам. Энтропия такого состояния равна



\[

S=k \ln W

\]



где $k$ - постоянная Больцмана. Наиболее вероятному состоянию соответствует максимум энтропии (при заданных $E=\sum_{i} E_{i} N_{i}$ и $N=\sum_{i} N_{i}$ ) распределения Бозе-Эйнштейна



\[

\bar{n}_{i}=\frac{\bar{N}_{i}}{G_{i}}=\left[\exp \left(\left(E_{i}-\mu\right) / k T\right)-1\right]^{-1},

\]



где $\mu$ - химич. потенциал, определяемый из условия, что сумма всех $n_{i}$ равна полному числу частиц в системе. Энтропия идеального газа, подчиняющегося Б.-Э. с., равна



\[

S=\sum_{i} G_{i}\left[\left(1+\bar{n}_{i}\right) \ln \left(1+\bar{n}_{i}\right)-\bar{n}_{i} \ln \bar{n}_{i}\right] .

\]



Тепловые колебания кристаллич. решетки твердого тела описываются как возбуждения совокупности осцилляторов, соответствующих ее нормальным колебаниям. Их можно рассматривать как идеальный газ квазичастиц фононов, подчиняющихся Б.-Э. с. На основании этого представления удается правильно описать поведение твердых тел при низких температурах. К важным приложениям Б.-Э. с. относится теория излучения черного тела, опирающаяся на представление о квантах электромагнитного поля - фотонах. Последние подчиняются Б.-Э. с.: в этом случае $\mu=0$, $E_{i}=h v(v-$ частота равновесного излучения в полости объема $V$ ). При этом распределение Бозе-Эйнштейна дает Планка закон излучения для спектрального распределения энергии излучения абсолютно черного тела.



Б. -Э. с. для системы взаимодействующих частиц основана на методе Гиббса для квантовых систем. Она может быть реализована, если известны квантовые уровни системы $E_{n}$ и может быть вычислена статистич. сумма



\[

Z=\sum_{n} \exp \left(E_{n} / k T\right),

\]



где суммирование ведется по всем квантовым уровням $\boldsymbol{n}$ системы для состояний, удовлетворяющих условиям квантовой симметрии. Последнее условие определяет тип квантовой статистики. Задача вычисления $Z$ не сводится к простой комбинаторной задаче и очень сложна, если взаимодействие между частицами не мало. Ее можно несколько упростить, если выразить гамильтониан системы в представлении вторичного квантования (в представлении чисел заполнения квантовых уровней) через операторы вторично-





\end{document}